\section*{Chapitre \ref{ch:iv2:risk-validation-and-regulation} --- Risque, validation et réglementation}
\markboth{\protect\omTocMark{}{Réponses~: Risque, validation et réglementation}}{\protect\omTocMark{}{Réponses~: Risque, validation et réglementation}}

\begin{solution}{iq:iv2:risk-validation-and-regulation:1}
Cela fait environ 2.33~millions de dollars pour un jour, et environ 7.36~millions de dollars pour dix jours.

La volatilité en dollars est de 1~million de dollars, et le quantile à 1~\% vaut 2.326 écarts-types. Pour dix jours, on multiplie par $\sqrt{10} = 3.162$.

\iqwrong{Dire «~on ne peut pas perdre plus de 2.33~millions de dollars~». La VaR est le seuil que la perte dépasse un jour sur cent. Elle ne dit rien de l'ampleur de la perte ce jour-là.}
\end{solution}

\begin{solution}{iq:iv2:risk-validation-and-regulation:2}
Cela fait environ 2.10~millions de dollars, et la diversification économise environ 0.11~million de dollars.

Les volatilités en dollars sont 0.5 et 0.45. Alors
\[
\sigma = \sqrt{0.25 + 0.2025 + 2 \times 0.8 \times 0.5 \times 0.45} = 0.901.
\]
Et $0.901 \times 2.326 = 2.10$. Additionner les montants autonomes donne $0.95 \times 2.326 = 2.21$.

\iqwrong{Multiplier la corrélation par les volatilités et oublier le facteur 2 du terme croisé. Il faut pondérer ce terme par les deux volatilités en dollars et le compter deux fois.}
\end{solution}

\begin{solution}{iq:iv2:risk-validation-and-regulation:3}
L'expected shortfall est la perte moyenne sachant que la perte dépasse la VaR. Pour une loi normale centrée réduite, elle vaut 2.338 à 97.5~\% et 2.665 à 99~\%. Ainsi, l'expected shortfall à 97.5~\% est presque égal à la VaR à 99~\% de 2.326.

La formule est $\mathrm{ES}_{\alpha} = \varphi(z)/(1 - \alpha)$, où $z$ est le quantile. Donc $\varphi(1.96)/0.025 = 2.338$.

C'est pourquoi le passage des régulateurs à l'expected shortfall à 97.5~\% maintient des fonds propres à peu près constants pour un portefeuille normal, et n'en exige davantage que si les queues sont plus épaisses.

\iqwrong{Croire que l'expected shortfall est nettement plus élevée à tous les niveaux. À 97.5~\%, sous une loi normale, elle dépasse à peine l'ancienne VaR à 99~\%.}
\end{solution}

\begin{solution}{iq:iv2:risk-validation-and-regulation:4}
La VaR à 99~\% vaut environ 5, et l'expected shortfall environ 6.6.

Un pour cent de 250 fait 2.5 observations. Prendre la troisième plus mauvaise observation comme VaR est la convention prudente, et elle donne 5. Interpoler entre la deuxième et la troisième donne 5.5.

L'expected shortfall est la moyenne des 2.5 pires observations~: $(8 + 6 + 0.5 \times 5)/2.5 = 6.6$.

\iqwrong{Traiter un seul nombre comme la réponse sans nommer la convention. Avec si peu de points, la convention seule fait varier la VaR de 10~\%.}
\end{solution}

\begin{solution}{iq:iv2:risk-validation-and-regulation:5}
Le modèle se trouve dans la zone jaune, et le multiplicateur passe de 3 à 3.65.

Le vert correspond à 0 à 4 dépassements, le jaune à 5 à 9 dépassements avec des compléments de 0.40, 0.50, 0.65, 0.75 et 0.85, et le rouge à 10 dépassements ou plus, avec 1.00.

\iqwrong{Dire qu'avec 2.5 dépassements attendus, le modèle a échoué. La zone jaune signifie «~cela peut aller dans les deux sens~», et les superviseurs examinent les causes des dépassements.}
\end{solution}

\begin{solution}{iq:iv2:risk-validation-and-regulation:6}
La demi-vie est d'environ 11.2 jours, et la nouvelle estimation est d'environ 1.086~\%.

La demi-vie résout $0.94^k = 0.5$, donc $k = \ln 0.5/\ln 0.94 = 11.2$.

La variance vaut $0.94 \times 0.0001 + 0.06 \times 0.0004 = 0.000118$, et sa racine carrée est 1.086~\%.

\iqwrong{Mettre à jour la volatilité plutôt que la variance~: $0.94 \times 1~\% + 0.06 \times 2~\% = 1.06~\%$. C'est un autre nombre, et ce n'est pas le modèle.}
\end{solution}

\begin{solution}{iq:iv2:risk-validation-and-regulation:7}
La volatilité totale est d'environ 26.9~\%, le $R^2$ d'environ 45~\%, et le portefeuille de 100 titres présente environ 18.1~\%.

La volatilité systématique vaut $1.2 \times 15~\% = 18~\%$. Le total vaut $\sqrt{0.18^2 + 0.20^2} = 0.269$. Le $R^2$ est $0.0324/0.0724 = 44.8~\%$.

Dans le portefeuille, la volatilité idiosyncratique tombe à $20~\%/\sqrt{100} = 2~\%$. Le total est donc $\sqrt{0.18^2 + 0.02^2} = 18.1~\%$. Le bêta reste inchangé.

\iqwrong{Croire que la diversification la fait tendre vers zéro. Elle n'élimine que la composante idiosyncratique. Le risque de marché ne se diversifie pas.}
\end{solution}

\begin{solution}{iq:iv2:risk-validation-and-regulation:8}
Le drawdown maximal vaut 50~\%, de 130 à 65.

On conserve le sommet courant. En chaque point, on calcule la baisse depuis ce sommet, et on retient la plus grande. Cela demande un seul parcours avec deux variables.

\iqwrong{Utiliser les points le plus haut et le plus bas de toute la série. Le creux doit venir après le sommet~; ici, cela tombe sur la même paire, mais un creux situé avant le sommet ne compte pas.}
\end{solution}

\begin{solution}{iq:iv2:risk-validation-and-regulation:9}
Cela fait environ 2.33~dollars par option. C'est trop élevé, car le gamma avantage une option longue.

Le mouvement vaut $2.326 \times \$2 = \$4.65$, et la moitié de ce montant est la perte en delta. Avec le gamma, $-0.5 \times 4.65 + 0.5 \times 0.02 \times 4.65^2 = -2.11$.

Pour une option courte, l'erreur va dans l'autre sens, ce qui est pire~: l'approximation delta-normale sous-estime alors la perte.

\iqwrong{Négliger le gamma au motif qu'il est petit. Sur un mouvement de 2.3 écarts-types, il représente 10~\% du chiffre ici, et pour les options courtes proches de l'échéance, il peut constituer l'essentiel du risque.}
\end{solution}

\begin{solution}{iq:iv2:risk-validation-and-regulation:10}
Cela fait environ 698\,000~dollars.

$\$50\,000 \times 6 \times 2.326 = \$697\,800$.

\iqwrong{Confondre la volatilité d'un taux exprimée en points de base avec sa volatilité en pourcentage du niveau. Une volatilité de taux cotée en pourcentage exige le niveau pour être convertie en points de base.}
\end{solution}

\begin{solution}{iq:iv2:risk-validation-and-regulation:11}
La VaR à 95~\% de chaque obligation vaut 0. Celle de la paire vaut 100. L'expected shortfall à 95~\% vaut 80 pour chacune et 103.2 pour la paire. La VaR n'est donc pas sous-additive~: combiner deux obligations l'a fait augmenter.

La perte de chaque obligation a une probabilité de 4~\%, ce qui reste dans la queue de 5~\%, si bien que son quantile à 5~\% vaut toujours zéro. Pour la paire, la probabilité d'au moins un défaut vaut $1 - 0.96^2 = 7.84~\%$, supérieure à 5~\%, donc la VaR vaut 100.

Pour l'expected shortfall, la queue de 5~\% de chaque obligation contient 4~\% à 100 et 1~\% à 0, soit $0.04 \times 100/0.05 = 80$. La queue de la paire contient 0.16~\% à 200 et 4.84~\% à 100, soit $(0.32 + 4.84)/0.05 = 103.2$, inférieur à $160$.

\iqwrong{Considérer la sous-additivité comme une simple question technique. Ici, une limite de VaR récompense la concentration~: on prend des risques de pertes rares qui se situent juste au-delà du quantile, et le chiffre reste à 0.}
\end{solution}

\begin{solution}{iq:iv2:risk-validation-and-regulation:12}
Oui, on rejette~: le rapport de vraisemblance vaut 5.50, au-dessus du seuil de 3.84, avec une $p$-valeur de 1.9~\%. La queue exacte vaut 1.37~\%. Avec 4 dépassements, le rapport vaut 0.77 et on ne rejette pas. Avec 0, il vaut 5.03, et on rejette dans l'autre sens.

\[
\mathrm{LR} = -2\ln\frac{(1 - p)^{n - x}p^{x}}{(1 - \hat p)^{n - x}\hat p^{x}},
\]
avec $\hat p = x/n$, à comparer à un $\chi^2_1$. La valeur exacte, $P(X \ge 7)$ pour une loi binomiale de paramètres $n = 250$ et $p = 1~\%$, sert de contrôle au test asymptotique.

Le cas sans aucun dépassement surprend. Sans aucune exception sur 250 jours, $\mathrm{LR} = -2 \times 250\ln(0.99) = 5.03$. La VaR est trop élevée. Cela coûte du capital, même si c'est une erreur plus prudente.

\iqwrong{Ne considérer qu'un seul côté. Le test est bilatéral, et un modèle trop prudent a lui aussi échoué. Pourtant, le feu tricolore ne sanctionne que le côté haut.}
\end{solution}

\begin{solution}{iq:iv2:risk-validation-and-regulation:13}
Les probabilités de transition valent 2.06~\% après un jour calme et 42.9~\% après une exception, le rapport de vraisemblance vaut 12.5, et l'indépendance est nettement rejetée, avec $p \approx 0.0004$.

$\pi_{01} = 5/243$, $\pi_{11} = 3/7$, et le taux agrégé vaut $8/250 = 3.2~\%$. Le rapport vaut
\[
-2\bigl[\ln L(\text{pooled}) - \ln L(\pi_{01}, \pi_{11})\bigr] = 12.5,
\]
à comparer au seuil de 3.84 d'un $\chi^2_1$.

Les exceptions groupées font que la VaR ne réagit pas aux changements de volatilité. Ainsi, quand un mauvais jour arrive, d'autres suivent, justement au moment où il faut que le chiffre soit juste.

\iqwrong{Valider le modèle parce qu'il paraît correct dans l'ensemble. Le test de comptage ne tient pas compte du moment où les exceptions surviennent, si bien qu'un modèle peut réussir ce test avec toutes les exceptions regroupées sur une seule semaine.}
\iqfollow{Que fait le test de couverture conditionnelle~?}{Il additionne les statistiques de Kupiec et d'indépendance, et compare la somme à un $\chi^2_2$.}
\end{solution}

\begin{solution}{iq:iv2:risk-validation-and-regulation:14}
De environ 20~\%~: la volatilité à dix jours vaut 3.79 fois la volatilité quotidienne, et non 3.16.

Pour un AR(1),
\[
\mathrm{Var}\Bigl(\sum_{i=1}^{m} r_i\Bigr) = \sigma^2\Bigl(m + 2\sum_{k=1}^{m-1}(m - k)\rho^k\Bigr).
\]
À $m = 10$ et $\rho = 0.2$, le crochet vaut 14.375, et sa racine carrée vaut 3.79. Le rapport vaut $3.79/3.16 = 1.20$.

En cas d'autocorrélation négative, la règle surestime, car des rendements qui reviennent à la moyenne s'accumulent moins. Et la formule redonne $\sqrt{m}$ pour $\rho = 0$.

Par ailleurs, la volatilité elle-même n'est pas constante, et la règle ignore aussi cet aspect.

\iqwrong{Appliquer $\sqrt{10}$ à tout. Les actifs peu liquides dont les prix sont lissés présentent en général une autocorrélation positive, et pour ceux-là, c'est précisément le mauvais sens.}
\end{solution}

\begin{solution}{iq:iv2:risk-validation-and-regulation:15}
La VaR vaut environ 2.65, et le shortfall environ 3.69, contre 2.33 et 2.67 pour la loi normale.

Une loi $t_\nu$ a une variance $\nu/(\nu - 2) = 2$. L'échelle de variance unitaire est donc $\sqrt{1/2}$, et $3.747 \times 0.7071 = 2.65$.

Pour le shortfall, la forme fermée est
\[
\mathrm{ES} = \frac{f_\nu(q)}{0.01}\cdot\frac{\nu + q^2}{\nu - 1}\cdot\sqrt{\frac{\nu - 2}{\nu}} = 3.69.
\]

À variance égale, la VaR est donc supérieure de 14~\%, et le shortfall de 39~\%. Les queues épaisses pénalisent davantage le shortfall.

\iqwrong{Oublier de ramener la variance à 1. Avec le 3.747 brut, on comparerait une loi de variance 2 à une loi de variance 1.}
\end{solution}

\begin{solution}{iq:iv2:risk-validation-and-regulation:16}
Elle vaut environ 0.037 écart-type, soit environ 1.6~\% de la VaR. Pour 0.5~\%, il faut environ 103\,000 trajectoires.

L'erreur type asymptotique d'un quantile vaut $\sqrt{p(1 - p)/n}/f(q)$. Ici, $\sqrt{0.0099/10\,000} = 0.000995$, et $f(2.326) = 0.0267$. On obtient donc $0.0373$, soit $0.0373/2.326 = 1.6~\%$.

Comme elle décroît en $\sqrt{n}$, passer de 1.6~\% à 0.5~\% exige $(1.6/0.5)^2 \approx 10.3$ fois plus de trajectoires.

Par ailleurs, le shortfall a une erreur type différente, et des queues plus épaisses demandent davantage de trajectoires pour la même précision.

\iqwrong{Utiliser l'erreur type d'une moyenne, $1/\sqrt{n}$. L'estimateur d'un quantile dépend de la densité en ce point, et la densité est faible dans la queue.}
\end{solution}

\begin{solution}{iq:iv2:risk-validation-and-regulation:17}
Les contributions valent 1.110 et 0.987, soit un total de 2.097.

La covariance de chacune avec le total, divisée par la volatilité du total, donne $(0.25 + 0.18)/0.901 = 0.477$ et $(0.2025 + 0.18)/0.901 = 0.424$. On multiplie par 2.326.

Elles s'additionnent grâce au \index{allocation d'Euler}théorème d'Euler. La volatilité et la VaR gaussienne sont homogènes de degré 1 en la taille des positions, donc la somme de la taille multipliée par le risque marginal égale le total.

Pour le vérifier, on augmente la première position de 1~\%, on recalcule la VaR, et la variation doit valoir 1~\% de sa contribution, soit 0.0111.

\iqwrong{Allouer selon la VaR isolée. Ces VaR s'additionnent à 2.21, et non au total, et elles ne disent rien de qui ajoute du risque à la marge. Une couverture recevrait une allocation positive.}
\end{solution}

\begin{solution}{iq:iv2:risk-validation-and-regulation:18}
La perte est de 10.5 millions de dollars, et les actions doivent encore baisser de 9~\% pour une perte de 15 millions de dollars.

Actions~: $50 \times -20~\% = -\$10$ millions. Taux~: $+\$40\,000 \times 100 = +\$4$ millions. Spreads~: $-\$30\,000 \times 150 = -\$4.5$ millions.

Les 4.5 millions restants sur 50 millions d'actions représentent 9~\%, soit $-29~\%$.

\iqwrong{Supposer que les taux jouent toujours en faveur du portefeuille. Qu'une baisse des taux s'accompagne d'une baisse des actions, c'est précisément l'hypothèse testée. Lors d'une alerte d'inflation, les deux jambes perdent ensemble, et le scénario de stress doit aussi retenir ce cas.}
\end{solution}

\begin{solution}{iq:iv2:risk-validation-and-regulation:19}
La volatilité vaut environ 23.9~\%, et les facteurs expliquent environ 81~\% de la variance.

Les expositions du portefeuille aux facteurs valent $b = B^\top w = (1.04, 0.19)$. La variance factorielle vaut
\[
1.04^2 \times 0.04 + 2 \times 1.04 \times 0.19 \times 0.006 + 0.19^2 \times 0.01 = 0.0460.
\]
La variance idiosyncratique vaut $0.25 \times 0.03 + 0.09 \times 0.02 + 0.04 \times 0.04 = 0.0109$. Le total vaut 0.0569.

\iqwrong{Utiliser les poids à la place des poids au carré pour la variance idiosyncratique. Ces parties sont indépendantes, donc elles s'additionnent avec les poids au carré, et c'est pour cela que la diversification fonctionne.}
\end{solution}

\begin{solution}{iq:iv2:risk-validation-and-regulation:20}
Elle vaut environ 208\,000 \$ pour une corrélation de 0.9, et environ 465\,000 \$ pour une corrélation de 0.5. Un changement de corrélation multiplie donc la VaR par 2.2 sans aucun changement de positions.

La volatilité du spread vaut $\$200\,000\sqrt{2 - 2\rho}$, et on multiplie par 2.326.

Une paire à faible VaR est donc longue de corrélation~: elle repose sur le maintien d'une corrélation élevée. Or cette corrélation est la plus faible au moment où l'on en a le plus besoin~: une crise peut dissocier une paire, et une OPA peut la briser en un après-midi. La corrélation récente rend une position la plus sûre juste avant qu'elle n'explose, et dimensionner selon ce chiffre revient à bâtir une position importante au mauvais moment.

\iqwrong{Considérer un brut de 20 millions de dollars et une VaR de 208\,000 \$ comme un petit risque. Regarder le brut et tester la corrélation en stress, car c'est le brut qu'il faut pouvoir sortir.}
\end{solution}

\begin{solution}{iq:iv2:risk-validation-and-regulation:21}
Les deux ont à peu près la même VaR~: 2.33 contre 2.57. Leurs expected shortfalls diffèrent d'un facteur dix~: 2.67 contre 26.4.

Pour B, la perte de 50 représente la moitié de la queue à 1~\%. La VaR est donc le point où la partie normale fournit les 0.5~\% restants~: $0.995\,\Phi(-q) = 0.005$, donc $q = 2.57$. L'expected shortfall vaut
\[
\frac{0.005 \times 50 + 0.995\,\varphi(q)}{0.01} = 26.4.
\]

Ainsi, l'expected shortfall tient compte de l'ampleur de la queue et il est sous-additif, ce qui explique pourquoi les régulateurs sont passés à cette mesure.

\textbf{Ce que l'on perd.} L'expected shortfall se backteste plus difficilement, car un seul nombre ne précise pas ce que l'on vérifie. De plus, il repose sur peu d'observations dans la queue profonde, si bien que son estimateur est moins stable. Il n'est pas élicitable à lui seul~: aucune fonction de score ne le retient comme meilleure prévision, alors qu'il l'est conjointement avec la VaR.

\iqwrong{Dire que l'expected shortfall résout le risque de queue. Il repose toujours sur les mêmes données, et si l'événement à 0.5~\% n'appartient pas à l'échantillon, le chiffre vaut 2.67.}
\end{solution}

\begin{solution}{iq:iv2:risk-validation-and-regulation:22}
Elle vaut environ 710\,000~\$ avec cinq tranches égales, et environ 485\,000~\$ si l'on peut tout vendre en une journée.

La VaR à un jour vaut $\$10 \text{ millions} \times 2~\% \times 2.326 = \$465\,000$.

En vendant par tranches égales, la position détenue le jour $k$ vaut $(6 - k)/5$. La variance totale vaut donc $\sum_{k=1}^{5}((6 - k)/5)^2 = 2.2$ variances journalières, soit un facteur de $\sqrt{2.2} = 1.483$. Cela donne 690\,000~\$. Le facteur général vaut $\sqrt{(n + 1)(2n + 1)/(6n)}$.

La demi-fourchette de 20 points de base sur les 10~millions de dollars ajoute 20\,000~\$.

\iqwrong{Utiliser $\sqrt{5}$. Cela suppose de garder toute la position pendant cinq jours, alors que l'on vend en cours de route, si bien que le facteur vaut 1.48 et non 2.24. Et la fourchette est un coût certain, elle s'ajoute donc directement.}
\iqfollow{Qu'est-ce qui manque~?}{L'impact de marché, qui croît avec le rythme de négociation. Vendre vite réduit ce risque mais coûte davantage en impact, et ce compromis est le problème d'exécution optimale.}
\end{solution}

\begin{solution}{iq:iv2:risk-validation-and-regulation:23}
Commencer par définir à quoi il sert, puis vérifier la théorie, le reconstruire de manière indépendante, tester ses limites, le confronter aux résultats, et n'approuver qu'à l'intérieur d'un champ d'utilisation énoncé. On refuserait si l'on ne pouvait pas reproduire ses chiffres ou si on l'utilise en dehors de ses hypothèses.

\begin{enumerate}
\item \textbf{Périmètre.} Ce qu'il valorise, ses utilisateurs, et à quel point le P\&L et le capital en dépendent. Cela fixe la profondeur de la revue.
\item \textbf{Théorie.} Ses hypothèses sont-elles adaptées au produit~? Un modèle à un facteur appliqué à un produit dont le paiement dépend de la corrélation de deux taux échoue ici, quoi que fasse le code.
\item \textbf{Réimplémentation indépendante.} La reconstruire à partir de la documentation et comparer les résultats. Cela fait apparaître ce que le développeur ne voit pas.
\item \textbf{Références et limites.} Cas à forme fermée, parité, volatilité nulle ou très forte, et échéance très courte. Vérifier la convergence.
\item \textbf{Entrées.} D'où viennent les données de marché, comment le modèle est calibré, et à quel point la calibration est stable d'un jour à l'autre.
\item \textbf{Analyse des résultats.} Explication du P\&L, simulation de couverture, et écart entre le prix du modèle et les prix observés sur le marché ainsi que les services de consensus.
\item \textbf{Conditions.} Limites d'usage, réserves pour l'incertitude, date de revue et suivi.
\end{enumerate}

\begin{dated}{2026-11}{Les recommandations prudentielles}\label{dat:iv2:risk-validation-and-regulation:mrm}
Aux États-Unis, la SR 26-2 d'avril 2026 a remplacé les recommandations antérieures sur le risque de modèle, la SR 11-7, et elle maintient la remise en cause effective (effective challenge) par des experts objectifs, avec l'analyse des résultats et le benchmarking dans la validation. Au Royaume-Uni, la SS1/23 de la PRA est en vigueur depuis mai 2024, avec cinq principes. Il faut consulter le texte en vigueur avant de le citer.
\end{dated}

\iqwrong{Considérer la validation comme une nouvelle exécution des tests du développeur. Ceux-ci montrent que le code fait ce que le développeur voulait. La validation demande si c'était bien la chose à vouloir.}
\iqfollow{Le desk affirme que la transaction doit être passée aujourd'hui, et la validation prend des semaines. Que faites-vous~?}{Proposer une approbation conditionnelle, avec une limite de position, une réserve et une date, et la consigner par écrit. Le pire résultat est un modèle mis en production sans que personne ait précisé ce qu'il peut faire.}
\end{solution}

\begin{solution}{iq:iv2:risk-validation-and-regulation:24}
Si le modèle est juste, il reste en zone verte avec une probabilité de 89.2~\%. À 2~\%, il reste en zone verte 43.9~\% du temps, et passe en zone rouge 3.0~\%. À 3~\%, zone verte 12.8~\%, et zone rouge 22.1~\%.

La zone verte correspond à au plus 4 exceptions, et la zone rouge à au moins 10. Loi binomiale avec $n = 250$. Donc $P(X \le 4 \mid p = 1~\%) = 0.892$. Pour $p = 2~\%$, $P(X \le 4) = 0.439$, et $P(X \ge 10) = 0.030$.

Un modèle erroné d'un facteur deux passe donc en zone verte près de la moitié du temps. Et un modèle juste passe en zone jaune 11~\% du temps. C'est pourquoi un comptage sur un an est un test faible. Pour distinguer nettement 1~\% de 2~\%, il faut beaucoup plus de données ou un test qui exploite davantage que des comptages.

\iqwrong{Considérer la zone verte comme une preuve que le modèle est juste. Une zone verte est compatible avec un taux d'exceptions doublé.}
\end{solution}

\begin{solution}{iq:iv2:risk-validation-and-regulation:25}
L'exigence de capital est de 109.5. La VaR de la veille devrait dépasser 36.5 pour compter. La VaR stressée ajoute 73, soit le double des 36.5 du reste.

$\max(12, 3.65 \times 10) = 36.5$. $\max(18, 3.65 \times 20) = 73$. Donc la somme vaut 109.5.

Comme le multiplicateur est supérieur à 3, c'est presque toujours la moyenne qui l'emporte. Pour le premier terme, il faudrait un saut à plus de 3.65 fois la moyenne récente.

La VaR stressée applique le portefeuille actuel à des données issues d'une période de tensions financières sur 12 mois. Elle ne baisse pas en périodes calmes, et l'exigence de capital est donc moins procyclique qu'avec la seule VaR.

\begin{dated}{2026-11}{Où en sont les règles}\label{dat:iv2:risk-validation-and-regulation:frtb}
La norme de janvier 2019 sur le risque de marché remplace cette mesure par l'expected shortfall à 97.5~\%. L'application varie selon les pays et les dates. À mi-2026, la Commission européenne avait adopté des amendements limités dans le temps, applicables à partir du 1\up{er} janvier 2027, et la PRA du Royaume-Uni a fixé des changements de modèles internes pour le 1\up{er} janvier 2028. Vérifier la situation actuelle avant de la citer.
\end{dated}

\iqwrong{Croire que l'exigence de capital suit la VaR de la veille. Un saut soudain sur un jour n'a presque aucun effet, si bien que c'est la moyenne qu'un desk doit piloter.}
\end{solution}

\begin{solution}{iq:iv2:risk-validation-and-regulation:26}
La VaR delta-normale est nulle. La VaR exacte à 99~\% vaut environ 6\,635~\$. L'approximation par une normale à moyenne et variance donne environ 4\,290~\$, soit environ 35~\% de moins.

La perte vaut $0.5 \times 500 \times x^2$ avec $x = 2z$, soit $1\,000 z^2$. Son point à 99~\% provient de $z^2$ considéré comme un $\chi^2_1$, dont le quantile à 99~\% vaut $6.635$. La VaR vaut donc 6\,635~\$. Les deux queues du titre comptent, et $z = 2.576$ est le point bilatéral à 1~\%.

L'approximation a pour moyenne 1\,000 et pour écart-type $1\,000\sqrt{2} = 1\,414$. Donc $1\,000 + 2.326 \times 1\,414 = 4\,290$.

Le P\&L est fortement asymétrique, si bien qu'une forme normale ne rend pas compte de la queue.

\iqwrong{Déclarer aucun risque puisque le delta est nul. Ce portefeuille n'a aucun risque du premier ordre, et tout son risque se situe au second ordre, que la méthode ignore.}
\end{solution}

\begin{solution}{iq:iv2:risk-validation-and-regulation:27}
Le quantile de Cornish-Fisher vaut environ $-3.30$. Pour une loi $t_6$, la formule donne environ $-3.03$, alors que la valeur vraie est $-2.57$. Elle s'écarte donc de 18~\%. C'est une série, et non une réponse sûre dans les queues épaisses.

\[
z_{\mathrm{CF}} = z + \frac{(z^2 - 1)S}{6} + \frac{(z^3 - 3z)K}{24} - \frac{(2z^3 - 5z)S^2}{36}.
\]
Avec $z = -2.326$, $S = -0.5$ et $K = 3$, cela donne $-3.30$.

La loi $t_6$ a pour variance $6/4$, donc son quantile à variance unitaire vaut $-3.143 \times 0.8165 = -2.57$. Son kurtosis d'excès vaut 3, et $S = 0$, si bien que la formule donne $-2.326 - 0.70 = -3.03$.

La même kurtosis peut donc correspondre à des queues différentes. Le développement n'est exact que pour de faibles écarts à la loi normale, et il peut même produire des quantiles non monotones lorsque l'asymétrie et la kurtosis sont plus élevées. C'est une correction utile pour de petits $S$ et $K$. Pour des queues épaisses, il faut ajuster une distribution ou utiliser la queue empirique.

\iqwrong{Supposer qu'une correction par les moments d'ordre supérieur est toujours plus proche de la vérité. Ici elle s'écarte davantage de la vérité que la normale simple, en ampleur~: 0.46 de trop loin contre 0.24 de trop près. Et l'asymétrie et la kurtosis empiriques sont elles-mêmes bruitées.}
\end{solution}

\begin{solution}{iq:iv2:risk-validation-and-regulation:28}
La moyenne vaut 0 si le modèle est juste, et environ $-1.16$ si la volatilité vraie vaut 1.2 fois celle du modèle. Un seuil à 5~\% vaut environ $-0.70$, par simulation, et il détecte le cas à 1.2 environ 77~\% du temps sur 250 jours.

Sous le modèle, l'espérance de chaque jour vaut $E[X I] = -\alpha\,\mathrm{ES}$, si bien que $Z$ a pour moyenne $1 - 1 = 0$.

Si la volatilité vraie vaut $s$, $E[X I] = -s\,\varphi(v/s)$, avec $v = 1.96$. Donc la moyenne de $Z$ vaut $1 - s\,\varphi(v/s)/(0.025\,\mathrm{ES})$, et pour $s = 1.2$ cela donne $-1.158$.

La statistique a une bonne propriété pour une banque~: comme elle utilise les ampleurs en plus des comptages, un expected shortfall trop faible apparaît même si la VaR est juste. Le seuil doit provenir de la simulation, et il dépend du modèle retenu, même s'il reste stable pour une gamme de queues.

\iqwrong{Penser qu'on ne peut pas du tout effectuer de backtest sur l'expected shortfall. Son absence d'élicitabilité exclut un type de comparaison, mais pas tous les tests. Ce qu'il faut, c'est la prévision propre du modèle, jour après jour.}
\end{solution}

\begin{solution}{iq:iv2:risk-validation-and-regulation:29}
C'est environ 12~\% par jour au départ, donc environ 2.2 dépassements sur 20 jours, contre 0.2 attendu. Le modèle à pondération exponentielle est à environ 3.6~\% le dixième jour.

Au début, le seuil reste à 2.326 anciennes volatilités, soit 1.163 nouvelles, et $\Phi(-1.163) = 12.2~\%$. Maintenu fixe pendant 20 jours, cela donnerait 2.45 dépassements. Mais les nouvelles journées entrent dans la fenêtre comme les pires, et chacune relève un peu le troisième pire. Dans une simulation, le taux baisse d'environ 12.7~\% à environ 9~\% sur les 20 jours, et le total vaut environ 2.2.

Pour le modèle exponentiel, après 10 jours, une part $1 - 0.94^{10} = 46~\%$ du poids porte sur le nouveau régime. La variance vaut donc $0.54 + 0.46 \times 4 = 2.38$ anciennes unités, et la volatilité 1.54 anciennes unités. Le seuil vaut $2.326 \times 1.544/2 = 1.80$ nouvelles volatilités, donc $\Phi(-1.80) = 3.6~\%$.

Les dépassements de la VaR historique se regroupent donc, et ils surviennent juste après les sauts de volatilité. Un modèle exponentiel rattrape son retard, même s'il part encore en arrière. Et la VaR du modèle historique conserve le chiffre élevé pendant un an après la fin de la tempête.

\iqwrong{Relever le niveau de confiance pour résoudre le problème. Un quantile plus élevé est un niveau plus élevé, mais il ne s'adapte pas plus vite.}
\end{solution}

\begin{solution}{iq:iv2:risk-validation-and-regulation:30}
La VaR vaut environ 3.90, et l'expected shortfall environ 5.62. Lorsque le paramètre de forme atteint 1, l'expected shortfall devient infini.

Avec $u = 2$, $\beta = 1$, $\xi = 0.2$, et un taux de dépassement de 5~\%,
\[
\mathrm{VaR}_p = u + \frac{\beta}{\xi}\Bigl(\Bigl(\frac{1 - p}{0.05}\Bigr)^{-\xi} - 1\Bigr) = 2 + 5(0.2^{-0.2} - 1) = 3.90.
\]
Et
\[
\mathrm{ES}_p = \frac{\mathrm{VaR}_p + \beta - \xi u}{1 - \xi} = \frac{3.90 + 1 - 0.4}{0.8} = 5.62.
\]

Le $1 - \xi$ au dénominateur en est la raison. Une forme égale ou supérieure à 1 signifie que la queue n'a pas de moyenne finie, donc qu'aucune perte moyenne n'existe à communiquer. Pour une forme nulle, la queue est exponentielle~: la VaR vaut $2 + \ln 5 = 3.61$, et l'expected shortfall la dépasse exactement de $\beta$.

\iqwrong{Choisir un seuil si élevé qu'il ne reste que peu de points, ou si bas que la forme de Pareto ne tient pas. Tracer l'estimation de $\xi$ en fonction du seuil, et retenir une portion stable.}
\end{solution}

\begin{solution}{iq:iv2:risk-validation-and-regulation:31}
C'est environ 0.13~\% sous la copule gaussienne, et environ 0.29~\% sous la $t_4$, contre 0.01~\% pour l'indépendance. La limite du coefficient de dépendance de queue vaut 0 pour la gaussienne et environ 0.25 pour la $t_4$.

Le chiffre gaussien est la normale bivariée en $(c, c)$, avec $c = -2.326$, et le chiffre de la loi $t$ provient d'une simulation. Ainsi, sachant que l'une franchit son point de 1~\%, l'autre suit avec une probabilité d'environ 13~\% ou 29~\%.

Le coefficient de dépendance de queue d'une copule $t$ vaut
\[
\lambda = 2\,t_{\nu + 1}\Bigl(-\sqrt{\frac{(\nu + 1)(1 - \rho)}{1 + \rho}}\Bigr) = 0.253.
\]
Pour la gaussienne, $\lambda = 0$ dès que $|\rho| < 1$~: à la limite, les extrêmes sont indépendants.

Une même corrélation peut donc signifier deux fois plus de risque de chute conjointe, et une copule gaussienne aux marges épaisses le manque encore.

\iqwrong{Penser que des marges à queues épaisses corrigent les queues conjointes. Elles modifient la queue propre de chaque actif, mais c'est la copule qui décide si elles chutent ensemble.}
\end{solution}

\begin{solution}{iq:iv2:risk-validation-and-regulation:32}
La mesure de capital vaut 16. Ce n'est pas $10\sqrt{12} = 34.6$, car la mise à l'échelle de chaque horizon ne s'applique qu'aux facteurs de cet horizon, et les facteurs ne subissent de choc que sur le temps supplémentaire.

\[
\mathrm{ES} = \sqrt{\mathrm{ES}_T(P)^2 + \sum_{j \ge 2}\Bigl(\mathrm{ES}_T(P, j)\sqrt{\frac{\mathrm{LH}_j - \mathrm{LH}_{j-1}}{T}}\Bigr)^{2}}.
\]
Les termes valent $10^2$, $(8\sqrt{1})^2 = 64$, $(5\sqrt{2})^2 = 50$, $(3\sqrt{2})^2 = 18$ et $(2\sqrt{6})^2 = 24$. Le total fait 256, soit 16.

S'il n'existe que des facteurs à 10 jours, les termes suivants s'annulent, et l'on obtient 10. Si chaque facteur est à 120 jours, tous valent 10, et la somme vaut $100(1 + 1 + 2 + 2 + 6) = 1\,200$, soit $10\sqrt{12} = 34.6$. Ce chiffre naïf est donc la borne supérieure, et la formule passe ce contrôle.

\iqwrong{Mettre à l'échelle l'ensemble du portefeuille sur son horizon le plus long. C'est la limite haute, et cela fait payer aux facteurs liquides un horizon d'illiquidité.}
\end{solution}

\begin{solution}{iq:iv2:risk-validation-and-regulation:33}
La volatilité et la VaR sont divisées par deux. On s'attendrait alors à environ 31 dépassements par an dans un test à 99~\%, contre 2.5. Un validateur le repère en refaisant à la main un petit calcul avant toute exécution sur les données.

Pour de petites variations, $x_t + x_{t-1} \approx 2x_{t-1}$, si bien que le ratio vaut la moitié du ratio prévu, et chaque rendement est divisé par deux. Chaque volatilité est donc divisée par deux, et la VaR avec elle. C'est exact au premier ordre, quelle que soit la volatilité.

Pour des rendements gaussiens, le seuil vaut $2.326/2 = 1.163$ écarts-types vrais. Le taux est donc $\Phi(-1.163) = 12.2~\%$, soit 31 jours sur 250.

\textbf{Comment le repérer tôt.} Calculer à la main un seul nombre. Comparer la volatilité en entrée à une source indépendante. Introduire une série dont la volatilité est connue et vérifier qu'elle la retrouve. Ne pas laisser le modèle de production vivre dans un tableur non relu, avec des étapes manuelles.

\iqwrong{Faire confiance au backtest pour le détecter. Si un changement de modèle intervient dans un marché calme, ou sur des positions couvertes, les dépassements peuvent mettre des mois à apparaître. Et un modèle qui abaisse la VaR est rarement remis en question par ceux qu'il arrange.}
\end{solution}

\begin{solution}{iq:iv2:risk-validation-and-regulation:34}
Il vaut 21 pour $V - C = 5$, et environ 5.34 pour $-20$. Le plancher du multiplicateur est de 5~\%, ce qui donne 0.7 lorsque RC est nul.

Pour $V - C = 5$, $\mathrm{RC} = 5$, et le multiplicateur est plafonné à 1~: $1.4 \times (5 + 10) = 21$.

Pour $-20$, $\mathrm{RC} = 0$, et le multiplicateur vaut $0.05 + 0.95\,e^{-20/19} = 0.382$. D'où $1.4 \times 3.82 = 5.34$.

Être hors de la monnaie ou surcollatéralisé réduit aussi l'exposition future, car il faudrait un mouvement pour revenir à une valeur positive. Lorsque $V - C \to -\infty$, le multiplicateur vaut $0.05$, et $1.4 \times 0.5 = 0.7$. Le plancher existe pour que le collatéral excédentaire ne ramène jamais l'exposition à zéro.

\iqwrong{Assimiler RC à l'exposition totale. L'exposition est aussi la variation potentielle sur le reste de la transaction, et c'est ce que représente l'add-on.}
\iqfollow{Pourquoi le 1.4~?}{C'est le facteur d'échelle prudentiel, $\alpha$, un nombre fixe fixé par la règle, qui couvre des éléments comme le wrong-way risk et l'erreur de modèle.}
\end{solution}

\begin{solution}{iq:iv2:risk-validation-and-regulation:35}
Le scénario le plus probable de ce type est $x_1 = -5.38$, $x_2 = -0.77$, et il se situe à environ 2.77 écarts-types de l'origine.

Sous des facteurs gaussiens, il faut maximiser la densité sachant la perte, ce qui revient à minimiser $x^\top\Sigma^{-1}x$ sous la contrainte $b^\top x = -10$. Avec un lagrangien,
\[
x^* = -\frac{10\,\Sigma b}{b^\top\Sigma b},
\]
avec $b = (2, -1)$, $\Sigma b = (7, 1)$ et $b^\top\Sigma b = 13$. D'où $x^* = -(10/13)(7, 1)$. La distance vaut $10/\sqrt{13} = 2.77$. Le P\&L a un écart-type de 3.61, donc 10 correspond à 2.77 de ces écarts-types.

\textbf{Le second facteur.} Vous êtes vendeur de ce facteur, donc vous gagneriez à une baisse, et une baisse de $x_2$ vous rapporte $+0.77$. Elle se produit de toute façon, car il est positivement corrélé au premier facteur, et une forte baisse de celui-ci l'entraîne. Le scénario est celui que la corrélation rend vraisemblable, et non celui qui fait le plus mal par unité.

Comparer avec un choc sur le premier facteur seul, $(-5, 0)$, qui perd également 10. Sa distance vaut 2.89, il est donc moins vraisemblable que $x^*$.

\iqwrong{Choquer un facteur à la fois et appeler reverse stress le pire résultat. L'objectif est de trouver la combinaison plausible, et elle n'est souvent pas celle que l'on aurait devinée.}
\end{solution}

\begin{solution}{iq:iv2:risk-validation-and-regulation:36}
Cette matrice n'est pas une matrice de corrélation. Sa plus petite valeur propre vaut $-0.547$. Avec les deux paires à 0.9, la corrélation (1, 3) doit appartenir à $[0.62, 1]$.

L'intervalle provient d'une condition sur le déterminant, ou de la corrélation partielle~:
\[
\rho_{13} \in \rho_{12}\rho_{23} \pm \sqrt{(1 - \rho_{12}^2)(1 - \rho_{23}^2)} = 0.81 \pm 0.19.
\]

Si on l'utilise malgré tout, il existe un portefeuille à variance négative~: on pondère selon le vecteur propre associé à la valeur propre négative, et la formule prend la racine carrée d'un nombre négatif, ou renvoie un risque nul, de sorte qu'un optimiseur s'y engouffre.

Pour la réparer, il faut modifier le moins possible~: ramener les valeurs propres négatives à une petite valeur positive, puis renormaliser à diagonale unitaire, ou résoudre un problème de plus proche corrélation. Il faut ensuite relire le résultat, car la réparation peut avoir déplacé les nombres que l'on voulait stresser.

\iqwrong{Penser que n'importe quel nombre entre $-1$ et 1 convient. Les bornes par paires ne suffisent pas, car les corrélations doivent être mutuellement cohérentes. Les tests de résistance qui fixent quelques paires à la main tombent constamment dans ce piège.}
\end{solution}

\begin{solution}{iq:iv2:risk-validation-and-regulation:37}
La moyenne reste proche de 2.5, mais la variance vaut environ 10.6 contre 2.475. Elle est donc environ 4.3 fois trop grande, et le compte se comporte comme si l'on n'avait que 58 observations indépendantes.

Les fenêtres se recouvrent sur neuf jours sur dix, donc une exception aujourd'hui tend à en rester une pendant plusieurs jours. La moyenne d'un comptage n'est pas affectée, puisque chaque fenêtre a une probabilité marginale de 1~\%~; mais les exceptions se regroupent en paquets pouvant atteindre une dizaine, ce qui fait monter la variance.

\textbf{Ce que cela change.} Kupiec suppose une loi binomiale. Avec la distribution réelle, plus large, on rejette un modèle correct bien trop souvent au seuil habituel. Et le test d'indépendance échoue par construction.

\textbf{Que faire.} Faire le backtest à un jour, ce que prescrivent les règles~: VaR à un jour pour le backtesting, VaR à dix jours pour le capital. Ou utiliser des fenêtres de dix jours sans recouvrement, bien que cela ne laisse que 25 observations par an. Ou simuler la loi sous l'hypothèse nulle du comptage avec recouvrement, et utiliser son seuil.

\iqwrong{Compter les fenêtres qui se recouvrent comme 250 observations de risque à dix jours. Il n'y en a que 25 indépendantes. Toute affirmation sur la fréquence des queues tirée de là doit rester prudente.}
\end{solution}

\begin{solution}{iq:iv2:risk-validation-and-regulation:38}
Le test de Kupiec passe, avec un rapport de 2.61 et une $p$-valeur de 10.6~\%. Le test d'indépendance échoue nettement, avec un rapport de 21.9 et une $p$-valeur d'environ 0.0003~\%. Le comptage est donc acceptable, mais le calendrier ne l'est pas. Le modèle réagit trop lentement.

\textbf{Chiffres.} Avec $\hat p = 9/500 = 1.8~\%$, Kupiec donne 2.61 contre un seuil de 3.84, et la queue exacte $P(X \ge 9) = 6.7~\%$. Les transitions donnent $\pi_{01} = 5/490 = 1.02~\%$ et $\pi_{11} = 4/9 = 44.4~\%$, soit un rapport de 21.9. Une exception rend la suivante plus de 40 fois plus probable.

\textbf{Ce que j'examinerais, dans l'ordre.}
\begin{enumerate}
\item \textbf{Quel P\&L~?} Faire le backtest sur le P\&L hypothétique, positions figées, sans frais ni négociation intrajournalière. Si les exceptions n'apparaissent que sur le P\&L réel, le problème vient peut-être des nouvelles transactions et non du modèle.
\item \textbf{Quand~?} Aligner les grappes sur les événements de marché. Si elles coïncident avec des sauts de volatilité, la fenêtre est trop lente~: tester une simulation historique ajustée à la volatilité ou une pondération exponentielle.
\item \textbf{Quelles positions~?} Attribuer chaque exception par facteur. Un seul desk, ou un facteur absent du modèle (une base, un skew de volatilité, un produit nouveau), signale un problème de couverture et non un modèle lent. La correction n'est pas la même.
\item \textbf{Données.} Les proxies périmés et les jours manquants lissent la volatilité et provoquent des dépassements lorsque les prix réels rattrapent les estimations.
\item \textbf{Positions mal dimensionnées.} Vérifier le gamma et le risque de saut à l'approche de l'échéance, car une approximation linéaire peut sembler juste jusqu'à un grand mouvement.
\item \textbf{Comparaison.} Faire tourner le modèle candidat sur le même historique~: le nouveau doit corriger le regroupement sans simplement gonfler le niveau.
\end{enumerate}

\textbf{Ce que je dirais.} Au desk~: le niveau est correct, mais le modèle est lent, et les limites qui en découlent sont trop lâches juste après un choc. Au régulateur~: exactement cela, avec les preuves, le plan de remédiation et sa date, ainsi que ce que nous faisons entre-temps, par exemple un supplément temporaire. Sur une horloge de 250 jours, une grappe peut aussi faire passer la banque en zone jaune~; il faut donc le signaler tôt plutôt qu'attendre que le multiplicateur change.

\iqwrong{Donner le 2.61 et s'en tenir là. Sur le seul comptage, ce modèle passe. Les modèles de risque échouent plus souvent dans le calendrier des exceptions que dans leur nombre total.}
\iqfollow{Pourquoi ne pas relever la VaR d'un facteur uniforme jusqu'à ce que le comptage soit correct~?}{Cela relève le niveau en permanence, ce qui coûte du capital, et laisse le regroupement en place. Les exceptions restent groupées, simplement un peu moins souvent.}
\end{solution}

\begin{solution}{iq:iv2:risk-validation-and-regulation:39}
L'approximation delta-normale donne 0. Le gamma seul donne une VaR à 99~\% d'environ 127\,000~\$. Gamma et vega ensemble donnent environ 174\,000~\$, avec un expected shortfall d'environ 223\,000~\$. Sous des mouvements de Student-$t$, la VaR vaut environ 252\,000~\$, et l'expected shortfall environ 525\,000~\$. Le scénario de stress perd environ 1.18~million de dollars. Je fixerais la limite sur le stress, et non sur la VaR.

\textbf{Le P\&L.} Avec $x$ le mouvement de l'action et $\Delta v$ la variation de volatilité,
\[
\mathrm{P\&L} = 22\,500 - 10\,000\,x^2 - 40\,000\,\Delta v.
\]
Le thêta est équilibré~: $\tfrac12 \times 20\,000 \times 1.5^2 = 22\,500$, donc le P\&L espéré est nul.

\textbf{Gamma seul.} $x^2 = 2.25 z^2$, dont le point à 99~\% vaut $2.25 \times 6.635$. La perte est donc $149\,285 - 22\,500 = \$126\,800$.

\textbf{Vega seul.} $40\,000 \times 2.326 = \$93\,000$. Les deux ajoutés donneraient 220\,000~\$. Ils ne sont pas parfaitement liés, et le vega ne nuit que dans une seule queue~; le chiffre combiné est donc de 174\,000~\$ par simulation de $\Delta v = -0.7 z + \sqrt{0.51}\,\varepsilon$. Une action qui baisse fait monter la volatilité~: les mouvements à la baisse sont touchés deux fois, et les mouvements à la hausse s'en compensent en partie.

\textbf{Queues épaisses.} On met les deux chocs à l'échelle d'une $t_4$ de variance unitaire, avec une échelle commune. La VaR augmente de 45~\%, et l'expected shortfall plus que double. Le gamma est un carré, il aggrave donc beaucoup plus les queues épaisses.

\textbf{Le stress.} $-\$10$ et $+5$ points~: $10\,000 \times 100 - 22\,500 + 200\,000 = \$1\,177\,500$. C'est près de sept fois la VaR normale. Un mouvement de $-\$10$ représente 6.7 écarts types. Sous une loi normale, cela n'arrive jamais. Dans un livre sur une seule action, une surprise de résultats ou une OPA peut le produire en un seul saut.

\textbf{La limite.} La fixer sur des scénarios~: une grille de mouvements de l'action et de la volatilité, avec une perte maximale sur la grille, par exemple $\pm 10~\%$ avec $\pm 5$ à $10$ points. La VaR ne voit pas cet écart, car c'est une quadratique qui croît avec le carré du mouvement. Il faut ensuite garder la VaR comme mesure quotidienne, et ajouter des limites sur le gamma et le vega par échéance.

\iqwrong{Donner le delta-normal pour un livre couvert en delta. Le chiffre qui affirme «~pas de risque~» est le premier dont chaque responsable des risques doit se méfier. Citer ce qui le met en défaut.}
\iqfollow{L'action saute de $+\$10$ et la volatilité baisse de 3 points. Et alors~?}{La perte vaut $1\,000\,000 - 22\,500 - 120\,000 = \$857\,500$. Un écart à la hausse coûte déjà beaucoup, donc le livre souffre dans les deux sens.}
\end{solution}

\begin{solution}{iq:iv2:risk-validation-and-regulation:40}
Le modèle est correct à la monnaie et dans la plage normale, il est précis au premier ordre avec une oscillation pair-impair, il est fortement décalé en valeur relative loin hors de la monnaie, et présente un vrai défaut à faible volatilité. Je l'approuverais sous conditions une fois ce défaut corrigé.

\textbf{Oscillation.} Les erreurs valent $-0.0399$, $+0.0344$, $-0.0200$, $+0.0174$, $-0.0100$, $+0.0087$. Le signe suit la parité de $N$, car le dernier nœud se trouve soit sur le prix d'exercice, soit de part et d'autre. La taille est divisée par deux quand $N$ double, donc l'ordre est $1/N$. \emph{Correction~:} faire la moyenne de deux pas adjacents. Pour $N = 100, 101$, la moyenne vaut 10.4493, soit une erreur de $-0.0013$, ce qui bat $N = 1\,000$ seul avec $-0.0020$. Ou lisser le dernier pas en utilisant la formule fermée à cet endroit.

\textbf{Loin hors de la monnaie.} Au prix d'exercice de 180, $N = 50$ est sous-estimé de $12.9~\%$ en valeur relative, et $N = 200$ de $3.6~\%$. L'erreur absolue est minime, mais un desk qui vend des ailes à bas prix se soucie du rapport. \emph{Correction~:} une tolérance relative par prix d'exercice, davantage de pas à cet endroit, ou une formule fermée lorsqu'elle existe. Il faut préciser la plage d'utilisation.

\textbf{Faible volatilité.} La probabilité de hausse vaut $p = (e^{r\Delta t} - d)/(u - d)$, et elle exige $\sigma > r\sqrt{\Delta t}$. Pour $\sigma = 0.1~\%$, $N = 500$, $\Delta t = 0.002$~: $r\sqrt{\Delta t} = 0.00224$, nettement au-dessus de $\sigma$. Donc $p > 1$, et les poids ne sont pas des probabilités. La récurrence rétrograde multiplie alors les erreurs à chaque pas, jusqu'à faire exploser le résultat. La limite devrait valoir $S - Ke^{-rT} = 4.88$. \emph{Correction~:} vérifier $0 \le p \le 1$ et refuser ou refaire la grille, ou utiliser plus de 2\,500 pas ici. Un arbre trinomial ou un réseau à dérive ajustée évite ce problème.

\textbf{Ce qu'il faut vérifier avant de signer.}
\begin{itemize}
\item Parité~: call moins put contre $S - Ke^{-rT}$ sur toute la grille. Avec ces données, l'accord vaut à 12 chiffres près.
\item Grecques~: vérifier les deltas et gammas obtenus par décalage, pour détecter le bruit. Le gamma d'un arbre peut osciller alors que le prix paraît correct.
\item Cas limites~: taux nul, échéance très courte ou très longue, volatilité très élevée, dividendes, option profondément dans la monnaie.
\item Exercice~: s'il s'agit d'options américaines, comparer avec une grille fine séparée, car il n'existe pas de formule fermée. Vérifier qu'un call américain sans dividende égale l'européen.
\item Actualisation~: un arbre affecté d'une erreur de taux reste proche de la bonne réponse. Tester un cas qui isole l'actualisation, comme le forward sans arbitrage.
\item Vitesse et mémoire, et comportement face aux données erronées.
\end{itemize}

\iqwrong{Comparer un seul prix à un seul nombre de pas, voir une concordance à deux décimales et signer. Un cas limite est ce qui révèle ce défaut, car aucune entrée typique ne le montre. Le pire échec est un grand nombre faux, sans aucune alerte levée.}
\end{solution}
